\rhead{EX101: EE Fundamentals Lab}
\phantomsection 
\section*{EE Fundamentals Lab (EX101)}
\label{course:EX101}

\noindent \small{\hyperref[main:scheme]{$\leftarrow$ Back to Scheme}} \vspace{0.5cm}

\noindent \textbf{Credits:} 2 (0-0-2)

\subsection*{Course Content}
\noindent\textbf{Lab 1: Verification of Circuit Laws (KCL and KVL)} \
Focuses on the fundamental conservation laws of electrical engineering. Students will construct a resistive network to verify Kirchhoff’s Current Law (node analysis) and Voltage Law (loop analysis) by measuring actual branch currents and voltage drops. \\

\vspace{0.2cm}

\noindent\textbf{Lab 2: Network Theorems (Thevenin, Norton, and Max Power)} \
Focuses on circuit simplification techniques. Students will experimentally determine the Thevenin equivalent voltage and resistance of a complex circuit and verify the Maximum Power Transfer theorem, which is crucial for impedance matching in signal transmission. \\

\vspace{0.2cm}

\noindent\textbf{Lab 3: Transient Response of RC Circuits (Time Constants)} \
Focuses on the timing behavior of circuits. Students will observe the charging and discharging curves of a Capacitor through a Resistor on an oscilloscope, calculating the Time Constant to understand propagation delays in digital logic. \\

\vspace{0.2cm}

\noindent\textbf{Lab 4: Passive Filters (Low-Pass and High-Pass)} \
Focuses on signal conditioning. Students will design and test RC filter circuits to understand how to attenuate high-frequency noise (Low-Pass) or block DC offsets (High-Pass), relating this to signal integrity in communication lines. \\

\vspace{0.2cm}

\noindent\textbf{Lab 5: Diode Characteristics and Logic Protection} \
Focuses on the P-N junction. Students will plot the V-I characteristics of a diode and implement a "Clipping Circuit," demonstrating how diodes are used to protect sensitive CPU pins from over-voltage spikes. \\

\vspace{0.2cm}

\noindent\textbf{Lab 6: The Transistor as a Digital Switch} \
Focuses on the core component of the CPU. Students will configure a BJT or MOSFET to switch a load (like an LED or Relay) ON and OFF, identifying the Cutoff (Logic 0) and Saturation (Logic 1) regions of operation. \\

\vspace{0.2cm}

\noindent\textbf{Lab 7: Linear Op-Amp Circuits (Math Operations)} \
Focuses on analog signal processing. Students will implement Inverting and Non-Inverting amplifiers and a Summing Amplifier using the 741 IC. \\

\vspace{0.2cm}

\noindent\textbf{Lab 8: Comparators and Digital-to-Analog Conversion} \
Focuses on the boundary between analog and digital. Students will build a Comparator circuit (1-bit ADC) and a simple 4-bit R-2R Ladder Network to generate analog voltages from digital inputs (DAC). \\

\vspace{0.2cm}

\noindent\textbf{Lab 9: Sensor Interfacing (LDR and Thermistors)} \
Focuses on reading the physical world. Students will interface resistive sensors (Light Dependent Resistor or Thermistor) using a Voltage Divider configuration to convert physical changes into measurable voltage changes. \\

\vspace{0.2cm}

\noindent\textbf{Lab 10: Pulse Width Modulation (PWM) and Motor Control} \
Focuses on actuation. Students will use a transistor driver to control the speed of a small DC motor or the brightness of an LED using a PWM signal (simulated via function generator), introducing the concept of duty cycle control.

\newpage
\rhead{Curriculum Schema}